\begin{proposition}[Three-by-five obstruction with explicit source hypotheses]
\label{B:three-by-five-finite-target-transfer}
Let $M$ be a finite monoid of order at most twenty-two with a
three-by-five rectangular minimal band, or the opposite shape.
Suppose that
\[
 H_*(BM;\Q)\cong H_*(\CP^2;\Q)
\]
as graded vector spaces, and that every integral degree-two class on
$BM$ restricts to zero along every map $BG\to BM$ from a finite group.
Then the square of a generator of $H^2(BM;\Z)\cong\Z$ is not primitive
in $H^4(BM;\Z)/\operatorname{torsion}$.
In particular the three-by-five lower-bound obstruction applies under
these hypotheses without assuming that $BM$ has finite CW type.
\end{proposition}

\begin{proof}
The minimal band makes the universal group trivial, so $H_1(BM;\Z)=0$.
The integral universal coefficient theorem and the rational homology
hypothesis therefore give $H^2(BM;\Z)=\Z$.  We allow torsion in
$H_2(BM;\Z)$ throughout the argument.

Apply the complete relation-space classification in
Theorem~\ref{B:three-by-five-targets}.  Its hull computations have three
types of output: explicit zero cup products, vanishing rational fourth
homology, and nonzero restrictions of actual integral degree-two
cocycles to eventual local cyclic groups.  The first two types pull
back along the actual monoid maps and require no finiteness assumption
on the source.  For the third type, an eventual local group in the
original source surjects onto the detected cyclic group.  Inflation of
its nonzero degree-two character is nonzero.  The assumed restriction
vanishing excludes exactly the same actors that Miller's theorem
excluded in the finite-CW proof.  Thus no finite-CW assumption on an
intermediate image or quotient is used.

The common repeated-column cases reduce to the preceding
three-by-three or three-by-four hull classifications.  Their zero
products, zero fourth homology and local-character tests apply for the
same reasons.  In the exceptional three-by-four parity case,
Proposition~\ref{B:even-hopf23} supplies a hull whose primitive
degree-two class has square twice an integral degree-four class.
The map to this hull is surjective on integral $H_2$: the row and
column augmentation quotients are surjective, and the balanced tensor
presentation adds relations.  Since the source has free $H_2$ rank
one and the target has $H_2=\Z$, the induced map on free parts is an
isomorphism.  Thus the target's primitive degree-two class pulls back
to a primitive source class, up to sign.  Its square is divisible by
two modulo torsion.  This step does not require torsion-free source
$H_2$.

The only remaining rank-one hull is orbit $164$.  Its forbidden-actor
test leaves the same allowed thirty-element hull and its $2186$
closed action images.  The profile quotient preserves integral $H_2$,
because it fixes the minimal band and the row and column actions.
There are $1385$ such images of order at most twenty-two.  A fresh
exact Smith computation for \emph{every} one of them gives
\[
\begin{array}{c|rrrrr}
 \operatorname{rank}H_2(-;\Z)&1&2&3&4&8\\\hline
 \text{number of images}&846&437&86&15&1.
\end{array}
\]
Every one of the $846$ rational-rank-one images has torsion-free
$H_2=\Z$.  Their indices agree exactly with the previously verified
small-image scope.  Consequently no image with an additional torsion
summand was omitted by the earlier integral filter.

For $842$ of these images, the existing independent rank certificates
give $H_4(-;\mathbb F_5)=0$, hence $H_4(-;\Q)=0$; all source
degree-two squares pulled back from such an image vanish rationally.
The remaining image indices are $1291,1296,1335,1376$, all of order
twenty-two.  Their first three normalized contracted bar chain groups
have dimensions $8,48,288$, and the two exact integer matrix ranks are
\[
 \operatorname{rank}_{\Q}d_1=7,\qquad
 \operatorname{rank}_{\Q}d_2=40.
\]
Thus each has $\dim H_3(-;\Q)=48-7-40=1$.
The profile quotient is surjective on $H_3(-;\Q)$ by
Lemma~\ref{B:band-quotient-H3-surjection}, contradicting
$H_3(BM;\Q)=0$.  This is a rational third-homology obstruction;
positive homology over a finite field alone would not suffice here.

The standalone checker
\path{code/b_orbit164_rational_small_images.py}
imports neither the action-image producer nor its bar implementation.
It reconstructs the actor multiplication, all $593910$ image closure
products, all $1385$ Smith presentations, and both exceptional
differentials.  It checks all $9216$ entries of the four products
$d_1d_2$, and computes their ranks over the integers.
The complete output is
\path{code/outputs/b_orbit164_rational_small_images_verified.json}.
This supplies the additional rational scope and torsion checks needed
for the present statement.
\end{proof}

\begin{remark}\label{B:three-by-five-transfer-scope}
The assumption on rational homology explicitly includes
$H_3(BM;\Q)=0$; bounded rational homology together with one-dimensional
$H^2$ and $H^4$ alone is not being substituted for it.  Only the
restrictions to the eventual local groups appearing in the cited hull
certificates are needed.  The stronger all-finite-group formulation
above is convenient when those restrictions are supplied by a map to
a finite simply connected target.  The proposition concerns band
sources.  For order at most twenty-two a three-by-five minimal ideal
with a nontrivial maximal group is impossible, since that ideal alone
would have at least thirty elements.
\end{remark}
